\chapter{Conclusion}
\label{chap:conclusion}

This thesis set out to close a narrow gap. The FAIR Jupyter reproducibility
pipeline can detect that a notebook failed because of a missing or
incompatible Python package, but it does not explain that failure and does
not try to repair it. \Cref{chap:related-work} showed that the pieces
needed to close this gap already existed separately: failure detection,
automated dependency inference, and LLM-based program repair. Among the
systems reviewed, none combined all three inside a reproducibility
pipeline, grounded in live package data and paired with a plain-language
explanation.

The repair layer built here does that. A deterministic classifier decides
scope and subtype before any language model is involved. An explanation
component turns a raw error message into a plain-language,
schema-validated account of what went wrong. Its scope is deliberately
wider than the repair scope: all 214 dependency-error records are eligible
for an explanation, not only the 200 that fall inside the pip-only repair
scope. A repair component proposes a fix using an open, locally deployable
model, checked field by field against live PyPI metadata and a small
curated compatibility registry, so a suggested package name or version is
never a guess. A fix-application component applies that proposal inside a
Docker environment rebuilt from the original recipe, re-runs the whole
notebook, and classifies the result. When that re-execution exposes a new
error, the layer classifies and explains it as well, and repairs it in one
bounded second round if it is repairable; explanation and repair stay
independent of each other throughout.

A result-logging component records one row per repair-decision entry,
including abstentions, in a SQLite benchmark table; in the final
evaluation every one of those entries was a repair-agent invocation. Exported to CSV and
run through the existing RML mapping, the final run's 235 rows became
3{,}478 RDF triples, in which each repair attempt is linked to the
notebook resource that the FAIR Jupyter knowledge graph uses for the
repaired notebook, validated in union with the existing notebook graph.
The
implementation is covered by 711 automated tests and was checked against
real PyPI, a real Ollama model, real Docker execution, real GitHub
repositories, and the real Morph-KGC mapping engine. The quantitative
evaluation used the 187 records reserved as the evaluation split; the 13
development records were used only for prompt design and pilots.

The three research questions of \Cref{sec:intro-rqs} are answered below,
each from the final frozen implementation run over that split.

\section*{RQ1: Explaining residual failures}

Gemma-2 9B completed 160 of the 187 original-failure explanations and all
50 explanations of errors newly exposed by a first repair; every one of
the 210 completed responses satisfied the explanation schema. The 27
missing records were timeouts and one service error of the local,
CPU-only model deployment, not malformed output. Fifteen participants,
with self-reported Python and Jupyter experience ranging from none to
regular use, each rated six of twelve frozen original-failure
explanations against a six-item questionnaire adapted from Hoffman et
al.~\cite{hoffman2023measures}. Of the 540 ratings, 88.9\% were Agree or
Strongly Agree, and every item reached a majority-agree median.
Understanding scored highest; satisfaction scored lowest, matching
free-text feedback that asked for more about the likely cause and less
restatement of the error message. The answer to RQ1 is therefore
positive within its scope: an open model produced structured explanations
for every failure it was able to answer, and the readers in this study
found them understandable, useful, and plausible. The ratings say nothing
about objective correctness, about a comparison with an alternative
method, or about the Round-2 and Qwen explanations, which no participant
rated.

\section*{RQ2: Resolving the targeted error and recovering the notebook}

The two halves of RQ2 have different answers, and they must be read
together. Every one of the 73 repair proposals the language model produced
was schema-valid and grounded in retrieved PyPI evidence, and 62 of the 73
FixApplicator attempts (84.9\%) removed the specific dependency error they
targeted. No notebook, however, re-executed cleanly
from start to finish within the two-round budget: full recovery is 0 of
187. The first figure is per attempt, the second per notebook, and they
are reported side by side throughout rather than combined into one
success rate.

Two factors explain the distance between them. The first is a deliberate
system behaviour: the repair agent's static import-to-distribution mapping
could not resolve the failing import in 162 of its 235 invocations, and it
abstained rather than guess a package name it could not ground. The second
is a property of the evaluated notebooks: many of them exposed further
dependency errors in sequence, so repairing one import frequently exposed
a second that had been invisible behind it; \Cref{sec:v-case-study} traces
a notebook that exposed a third dependency error after two successful
targeted repairs. The two-round cap then determines how far along such a
chain the system can follow. The answer to RQ2 is that the bounded
pipeline resolves the
targeted error in most of the attempts it makes, that it makes an attempt
for fewer than a third of its invocations, and that resolving one
dependency error did not restore complete execution within the two-round
budget in this evaluation.

\section*{RQ3: Sensitivity to the choice of model}

The same 187-notebook evaluation was repeated with Gemma-2 9B replaced by
Qwen3.6-35B-A3B-MLX-8bit. The pipeline, prompts, schemas, retrieval,
validation, evaluation records, and repair policy were held fixed; the
model, its provider, and the calibrated Qwen output-token budget differed.
Every aggregate repair figure was identical: the same 126 Round-1
abstentions, the same 73 grounded proposals, the same 62 targeted errors
resolved, and zero full recoveries for both. The two runs were not
identical in every decision. On 13 notebooks, all failing on the same
removed NumPy symbol, both models pinned NumPy from the same grounded
candidate list, Gemma to 1.26.3 and Qwen to 1.26.4. Both versions
satisfied the retrieved compatibility constraint, both installed, and the
notebooks then failed on the same next error under both models. Model
choice produced a small difference in version selection for these 13
repairs, but it did not change any measured downstream repair outcome.
Qwen also completed all 237 explanation records under its serving
environment, whereas the local Gemma/Ollama run lost 27 to runtime
timeouts. Because the model and the serving setup changed together, this
observation does not isolate the model as the cause. The answer
to RQ3 is that, under this pipeline, structured repair decisions were
nearly insensitive to the choice between these two models and the measured
repair outcomes were unaffected, while explanation completion differed
between the two runs. The retrieval and grounding layer bounds a valid
proposal to the resolved distribution and a small, PyPI-confirmed
candidate set, so what remains for a model to decide is a choice among a
few compatible versions. Neither model is shown to be generally better
than the other, and other models could behave differently.

\section*{Closing remark}

Taken together, the answer to the question this thesis opened with is a
qualified yes. A dependency-related notebook failure can be explained in
language a non-specialist reader finds understandable. The results show
that specific dependency errors can often be repaired automatically with
grounded, deterministically validated proposals from an open model, even
though complete notebook recovery remained out of reach within the
two-round budget. What this system cannot yet do is carry a real research
notebook all the way back to a clean run, because many notebooks in the
evaluation exposed additional dependency errors after an initial repair
and because resolving an import name to a package remains the binding
constraint. Two limits of the
evaluation itself also stand: the repositories were re-executed at their
default-branch state rather than at the originally failing commit, and the
local model service did not answer every explanation request in time.
\Cref{chap:future-work} sets out what would have to change for that, and
\Cref{sec:v-threats} states the boundaries within which every result above
should be read.
